\documentclass{ctexbook}
\usepackage{graphicx}
\usepackage{booktabs}
\usepackage{longtable}
\usepackage{array}
\usepackage{makecell}
\usepackage{amsmath}
\usepackage{natbib}
\title{基于深度学习的电梯边缘AI安全视觉关键技术研究与实现}
\author{叶伟键}
\date{}
\begin{document}
\maketitle
\begin{center}
Research and Implementation of Key Technologies for Elevator Edge AI Safety Vision Based on Deep Learning\\
学院：计算机科学与工程学院\\
专业：计算机科学与技术全英创新班\\
学生姓名：叶伟键\\
学生学号：202230430217\\
指导教师：钟竞辉
\end{center}

\chapter*{摘要}
电梯轿厢空间封闭、使用频繁，镜面反光、近距离遮挡和目标尺度突变经常同时出现。安全视觉系统既要识别轿厢内人员，也要判断电动车入梯风险；训练阶段得到的检测精度，迁移到边缘设备后未必能够原样保持。围绕这一问题，本文以 YOLOv8n 为检测模型，以搭载华为海思 Hi3516DV500 视觉 SoC 的 HongOU PI 开发板为部署平台，研究电梯轿厢中行人与电动车的双目标检测、模型迁移和部署一致性评价。

研究中将电梯安全视觉界定为 person 与 ebike 双目标检测任务，并结合轿厢成像特点讨论其安全需求。在模型迁移部分，训练权重、ONNX 导出、OM 编译和板端推理被放入同一条检查链路，输入、类别、输出、后处理和运行状态分别作为需要核对的约束。评价部分则不只看平均精度，还引入分类别指标、阈值敏感性、NMS 参数影响、后处理策略消融、分段计时和代表性视频片段，用来观察部署后的模型行为。

在 720 张验证图像上，板端总体 mAP@0.5 为 0.910，其中行人类别为 0.869，电动车类别为 0.952；相较训练阶段的 0.988，板端结果下降约 0.078，但主要检测能力仍被保留下来。置信度阈值在 0.20--0.40 范围内变化时，总体 F1 保持在 0.915--0.922 区间。分段计时中，单帧模型执行阶段耗时约 8.2 ms，能够支撑边缘侧实时处理。

\noindent\textbf{关键词：}电梯安全视觉；边缘智能；双目标检测；部署一致性；轿厢人车识别

\chapter*{Abstract}
Elevator cabins are enclosed and frequently used spaces, where mirror reflections, close-range occlusions, and rapid scale changes often occur together. A safety vision system in this setting has to detect passengers while also recognizing e-bikes entering the cabin. Accuracy measured during training, however, does not automatically carry over to an edge device. This thesis studies that gap with YOLOv8n as the detector and the HongOU PI board, based on the Huawei HiSilicon Hi3516DV500 vision SoC, as the deployment platform.

The task is formulated as dual-target detection for person and ebike, with the imaging conditions of elevator cabins used to explain the safety requirements. For migration, the trained weights, ONNX export, OM compilation, and board-side inference are treated as one inspection chain. Input, category, output, post-processing, and runtime constraints are checked separately. The evaluation then combines per-class metrics, threshold sensitivity, NMS parameter changes, post-processing ablation, segmented timing, and representative video clips, so that deployment behavior can be examined beyond a single average accuracy number.

On the 720-image validation set, the board-side system obtains 0.910 overall mAP@0.5, with 0.869 for person and 0.952 for ebike. The result is about 0.078 lower than the training-stage mAP@0.5 of 0.988, but the main detection capability is still retained after deployment. When the confidence threshold varies from 0.20 to 0.40, the overall F1 stays within 0.915--0.922. Segmented timing gives a per-frame model execution time of about 8.2 ms, which is compatible with edge-side real-time processing.

\noindent\textbf{Keywords: }Elevator safety vision; Edge intelligence; Dual-target detection; Deployment consistency; Cabin person-vehicle recognition

\chapter*{目录}
\begin{longtable}{p{0.82\textwidth}r}
摘 要 & I\\
Abstract & II\\
第一章 绪论 & 1\\
\quad 1.1 研究背景与问题挑战 & 1\\
\quad 1.2 国内外研究现状 & 2\\
\quad 1.3 研究目标与主要工作 & 4\\
\quad 1.4 研究内容组织逻辑与论文结构 & 7\\
第二章 技术基础与相关工作 & 8\\
\quad 2.1 目标检测基础 & 8\\
\quad 2.2 边缘部署技术基础 & 12\\
\quad 2.3 相邻视觉任务与本文重点 & 14\\
\quad 2.4 相关工作小结 & 15\\
第三章 YOLO 类双目标检测模型的边缘兼容迁移方法 & 17\\
\quad 3.1 总体技术路线 & 17\\
\quad 3.2 YOLOv8n 场景化训练与验证性能 & 18\\
\quad 3.3 迁移约束与一致性检查 & 19\\
\quad 3.4 一致性评价方法 & 22\\
\quad 3.5 本章小结 & 23\\
第四章 数据集、场景复杂性与评测协议 & 24\\
\quad 4.1 数据集与任务定义 & 24\\
\quad 4.2 评测协议设计 & 25\\
\quad 4.3 数据可追溯性与实验设计 & 26\\
\quad 4.4 本章小结 & 27\\
第五章 实验结果与分析 & 28\\
\quad 5.1 实验设置与评价范围 & 28\\
\quad 5.2 板端检测能力验证 & 29\\
\quad 5.3 部署参数敏感性分析 & 33\\
\quad 5.4 运行时评价 & 36\\
\quad 5.5 综合讨论 & 38\\
\quad 5.6 本章小结 & 40\\
结论 & 42\\
参考文献 & 43\\
致谢 & 49\\
\end{longtable}
\chapter{绪论}
\label{cha:introduction}

\section{研究背景与问题挑战}

电梯轿厢是城市建筑中使用频率很高的公共微空间。它的空间尺度小，人员进出快，摄像机视角又常被固定在顶部或角落，因此同一画面里容易同时出现近距离遮挡、镜面反光、金属门框高亮和透视畸变。通用视觉模型在道路、走廊或开放数据集上形成的经验，放到这类轿厢画面中并不总是适用。更麻烦的是，电梯本身属于安全敏感场景，系统既要看得见目标，也要在端侧稳定运行；持续漏检、误检或迁移后的行为漂移，都会削弱风险治理效果。

在现实应用层面，电梯安全视觉需要同时关注两类目标：行人和电动车。行人检测服务于轿厢人员计数、拥挤度感知和异常行为识别；电动车检测服务于入梯风险治理和安全预警。与普通静态物体检测任务不同，电梯场景中人与车常常同时出现、局部遮挡、仅露出车头或车把、人员身体压住车身、门口区域边框干扰等复杂情况。围绕行人与电动车建立统一的双目标边缘视觉方案，是比单类检测更贴近真实部署需求的技术路径。

边缘计算把推理过程尽量放到数据源附近，这一点与电梯安全视觉的需求较为契合\cite{shi2016edge}。原始视频不必长期上传云端，隐私和带宽压力会小一些；预警链路也更短，更适合电梯这种对时延敏感的场景。与此同时，设备端可以留下本地运行记录，便于回看模型行为。不过，边缘部署不是把模型格式换一下就结束了。模型轻量化、量化与编译、后处理阈值、跨格式输出解释，都会影响最终结果。

基于上述背景，本文聚焦电梯轿厢中的行人与电动车双目标检测，讨论场景化训练后的 YOLO 类模型如何在华为海思/昇腾边缘 AI 平台及工具链约束下完成兼容迁移。实验平台选用搭载华为海思 Hi3516DV500 视觉 SoC 的 HongOU PI 开发板，重点检查类别、边框、阈值和运行时行为在设备端是否仍能保持一致。


与普通目标检测相比，电梯轿厢边缘安全视觉存在更强的约束耦合。本文将这些约束概括为以下五类。

\begin{enumerate}
    \item 视角与尺度约束。电梯摄像机通常安装在轿厢顶部或角落，拍摄距离近、视角固定但透视变化明显。人在近景和门口区域会出现边框被截断、下半身出框或头肩放大的现象，电动车则常以局部结构出现。
    \item 反光与遮挡约束。金属门板、镜面装饰和玻璃表面会产生高亮或虚像，而人与电动车共同出现时又会互相遮挡，检测器的类别判断和边框定位会因此变得不稳定。
    \item 边缘部署约束。边缘设备受算力、功耗、存储和算子支持限制，模型结构不能无限增大，训练阶段形成的模型也不一定能无损迁移到部署阶段。
    \item 运行时稳定性约束。安全场景更关心连续视频中的表现，而不是孤立图像上的最好指标。某一帧检测正确，并不能保证系统在连续数秒内不会闪烁或类别跳变。
    \item 结论可复查约束。论文中的实验结论需要回溯到验证集、指标摘要和设备端运行结果，不能只停留在经验判断。
\end{enumerate}

图 \ref{fig:thesis_loop} 给出了本文采用的整体研究流程。该流程的关键不在于追求抽象的最佳模型，而在于把电梯场景中的难例分析、模型迁移与边缘一致性诊断统一到同一套评价框架中。

\begin{figure}[htbp]
    \centering
    \includegraphics[width=0.92\textwidth]{image/fig1-research-loop.png}
    \caption{本文采用的电梯边缘安全视觉研究流程}
    \label{fig:thesis_loop}
\end{figure}

\section{国内外研究现状}

\subsection{边缘智能与轻量部署方向}

边缘计算的基本思路，是把计算过程尽量放回数据产生的位置附近，以减少传输等待、带宽占用和隐私暴露\cite{shi2016edge,chen2024edgeai}。电梯视频正好具有这几类约束：画面涉及乘客隐私，预警又要求尽快响应，若全部交给云端，网络抖动会直接进入报警链路。围绕端侧运行，模型压缩研究已经给出多种工具。Jacob 等的整数推理友好量化说明，低位宽计算并不必然破坏神经网络性能\cite{jacob2018quantization}；Deep Compression 从剪枝、量化和编码入手压低模型存储与计算成本\cite{han2016deep}；知识蒸馏则提供了把大模型判别能力转移到轻量模型上的常见做法\cite{hinton2015distilling}。

不过，模型真的进入边缘设备后，问题并不止于“文件能不能编译”。训练框架里的浮点权重经过导出、量化或编译后，边框坐标、类别分数排序以及阈值响应都可能发生偏移。离线基准里，这类偏移有时只表现为几个小数点的变化；放到电梯安全视觉中，却可能变成连续画面里的误报、闪烁或漏检。因此，本文关注的不是单次编译成功，而是训练格式、离线模型和板端后处理之间的行为对照\cite{stacker2021deployment,chen2021aqd}。

\subsection{轻量目标检测与实时视觉方向}

实时目标检测已经在工业视觉中形成较成熟的工程路线。YOLO 系列采用端到端预测，结构相对简洁，在不少真实场景中都能直接落地\cite{redmon2017yolo9000}。不同计算预算下，检测网络通常需要在精度和计算量之间取舍：EfficientDet 通过跨尺度特征融合和复合缩放策略取得较好的折中\cite{tan2020efficientdet}；MobileNetV3、EfficientNet 分别从移动端架构搜索和统一缩放规律角度提供了轻量骨干设计参考\cite{howard2019mobilenetv3,tan2019efficientnet}。EdgeYOLO 等工作则把边缘实时检测作为目标，更关注算子友好性、端侧吞吐和推理效率\cite{zhou2023edgeyolo}。

围绕 YOLO 类模型的板端部署，已有研究也开始从算法精度扩展到嵌入式硬件实测。Feng 等在 Jetson Nano、Jetson Xavier NX、Raspberry Pi 4B 与 Intel NCS2 等平台上比较 YOLO 模型的推理流程、帧率和资源占用差异，说明同一检测器在不同边缘设备上会受到计算单元、运行时和模型格式的共同影响\cite{feng2022yolobenchmark}。YOLOBench 进一步在 x86 CPU、ARM CPU、NVIDIA GPU 和 NPU 等嵌入式硬件上评测大量 YOLO 模型，强调模型选择需要结合目标硬件的精度--时延权衡，而不能只依据通用数据集上的离线精度\cite{lazarevich2023yolobench}。这类工作表明，YOLO 边缘部署研究已经从“能否运行”转向“目标硬件上是否具备稳定、可解释的精度和效率表现”。

与上述通用嵌入式平台上的 YOLO 部署基准相比，电梯场景中的双目标检测更强调轻量性、稳定性和板端输出口径一致性。一方面，模型需要识别行人与电动车两个类别，类别数较少，理论上适合轻量模型；另一方面，轿厢内遮挡、反光和近距离尺度变化又会放大轻量模型的表达不足。在前期部署尝试中，YOLOv5 模型在 Hi3516DV500 的 NPU 上占用率可达约 80\%，资源余量有限，后续为追求更好的轻量化效果选择了 YOLOv8n。本文围绕该模型开展部署迁移与一致性评价。

\subsection{困难样本分析与数据增强方向}

当目标检测任务主要受长尾样本和困难样本影响时，单纯扩大数据量并不一定带来稳定收益。OHEM 等在线难例挖掘方法的思路，是让模型把更多训练权重放在当前最容易出错的样本上，从而改善关键失败模式\cite{shrivastava2016ohem}。目标检测中的定向数据增强也有类似作用。Scale-Aware Automatic Augmentation 等研究说明，围绕目标尺度、位置和场景结构设计增强策略，通常比无差别增强更贴近检测误差来源\cite{chen2021scaleaware}。

电梯轿厢中的困难样本具有明显场景规律：反光区域集中在门板和镜面结构，遮挡多发生在人车重叠区域，电动车局部可见常出现在门口或边缘位置。因此，本文将这些因素作为场景复杂性解释工具，用于说明指标变化与实际轿厢结构之间的关系，并将研究重点收束在双目标检测、板端兼容迁移和部署一致性评价上。

\subsection{相关场景研究与当前不足}

与通用监控或道路交通相比，直接面向电梯视觉的公开研究数量并不多。已有工作多聚焦于人数统计、乘客计数或局部行为识别，例如基于 MobileNet-SSD 的电梯人数统计系统研究表明，在电梯场景中采用边缘侧轻量模型是具备现实可行性的\cite{lin2023elevatorcounting}。另一些研究则尝试将轻量模型用于电梯跌倒检测，说明电梯作为独立应用场景具有实际研究价值。

不过，现有研究通常只覆盖单一任务或单一部署环节，较少同时讨论电梯轿厢中的行人与电动车双目标检测、场景复杂性解释以及训练阶段到边缘部署阶段的一致性问题。本文正是在这一不足上展开：以双目标检测为主线，以场景复杂性分析和部署一致性评价为两条支撑线，形成面向电梯安全视觉的研究框架。


\section{研究目标与主要工作}

\subsection{问题定义与核心思路}

从系统角度看，电梯边缘安全视觉不是单独训练一个检测器那么简单。数据来自封闭轿厢，模型要经过导出和编译，还要在开发板应用层完成阈值筛选与 NMS。训练阶段的高精度不能自动说明板端输出仍然可信；反过来，程序能跑起来，也不等于类别、边框和阈值含义没有走样。本文要回答的核心问题是：基于 YOLO 的 \textit{person + ebike} 双目标模型，在 HongOU PI（Hi3516DV500）这样的平台约束下，怎样完成兼容迁移，并用实测指标说明板端输出仍可复查。

这条问题链从轿厢场景本身出发。人车共现、遮挡反光和局部可见，使安全管理需求落到行人与电动车两个类别上；而模型迁移不是简单配置环境，输入尺寸、类别顺序、输出张量和后处理规则都需要逐项核对。本文把 PyTorch 权重、ONNX 中间表示、OM 离线模型以及板端应用后处理放在同一链路中检查。结果分析也采用组合证据：全量验证集说明总体能力，分类别指标暴露行人与电动车差异，阈值、NMS、后处理和分段计时则用来解释板端运行口径。

因此，本文把三个环节连在一起讨论：电梯安全需求被界定为双目标检测任务，模型进入国产边缘平台的过程被写入方法链路，一组可复查实验则用于观察设备端行为。这样处理后，检测精度不再是孤立数字，部署过程也不只是工程背景，而是判断系统能否用于电梯安全视觉的必要证据。

研究动机也来自这个落地过程。电梯安全视觉对漏检和误检都敏感，迁移中的小幅数值变化可能被放大为报警差异。YOLO 检测器虽然部署路径清晰，但输出张量、类别顺序和后处理参数一旦解释错误，板端结果就会偏离训练阶段预期。HongOU PI（Hi3516DV500）开发板采用海思视觉 SoC、专用 NPU 和 OM 离线模型格式，涉及 ONNX 导出、ATC 编译、板端运行时调用和应用层后处理，适合作为本文检查同类目标平台迁移约束的实验载体。

基于这些考虑，论文同时看两件事：模型能否识别行人与电动车，以及模型进入边缘设备后，输出语义和运行口径是否还能对得上训练阶段。只有把这两件事放到同一证据链里，电梯边缘安全视觉系统的可用性才说得清楚。

\subsection{研究目标与主要内容}

本文的研究目标是：面向电梯轿厢中的行人与电动车，完成 YOLO 类双目标检测模型的场景化训练、边缘侧迁移和板端实测分析。具体来说，论文整理从训练权重到 OM 模型再到 HongOU PI（Hi3516DV500）应用输出的检查流程，并用统一实验协议观察检测精度、阈值响应和运行时口径。轿厢中的反光、遮挡和局部可见目标不作为单独创新点展开，而是用于解释误差来源。围绕这一目标，本文完成以下工作。

\begin{enumerate}
    \item 建立电梯轿厢双目标检测任务定义，明确行人与电动车两个核心类别，其中行人检测服务于人员计数与异常感知，电动车检测服务于入梯风险识别。
    \item 总结电梯场景中的典型困难因素，用于解释遮挡、局部可见和光照变化对检测输出的影响。
    \item 分析 YOLO 类检测器在双目标任务中的结构适配方式，说明输入尺寸、类别顺序、检测头输出和 NMS 后处理为什么会影响板端结果。
    \item 梳理从 YOLO 类训练权重、ONNX、OM 到边缘设备执行的迁移链路，检查类别、边框、阈值和运行记录是否能与训练阶段对上。
    \item 基于统一验证集、参数敏感性、后处理消融、代表性片段和分段计时实验，报告总体指标、分类别指标、部署参数影响和边缘设备运行表现。
\end{enumerate}


\subsection{研究难点分析}

本文研究难点并不集中在单一环节。任务层面的困难来自轿厢内行人与电动车经常同时出现，并伴随遮挡、反光、局部可见和尺度变化等问题；模型需要在较小类别集合中保持稳定检测，而不是简单识别静态物体。

模型层面的困难在于，YOLO 类检测器虽然具有较好的实时性，但检测头输出、坐标解码、置信度筛选和 NMS 后处理共同决定最终结果。训练框架脚本通常自动完成这些步骤，而板端应用程序需要显式实现或调用相应逻辑。只要其中一个环节解释不一致，最终输出就可能变化。

部署层面的困难则更接近工程现场。目标平台需要把训练阶段模型转换为可执行的离线模型，并在板端运行时中解释输入、输出和后处理结果。作为本文的实验平台，HongOU PI（Hi3516DV500）开发板具有面向国产边缘 AI 的专用编译和运行时体系，转换过程中的算子适配、预处理配置、数据类型变化和输出读取方式，都会影响设备端行为。因此，部署链路必须进入论文方法本身，而不能只作为实验环境说明。

\subsection{本文的主要工作}

本文工作集中在任务建模、迁移检查和部署评价三个方面。

\begin{enumerate}
    \item \textbf{面向电梯轿厢的双目标检测任务建模。} 本文将电梯安全视觉需求建模为行人与电动车双目标检测问题，行人检测服务于人员计数与拥挤感知，电动车检测服务于入梯风险治理，并用场景复杂性因素解释遮挡、反光和局部可见等因素对检测输出的影响。
    \item \textbf{YOLO 类模型的边缘兼容迁移分析框架。} 本文面向海思/昇腾边缘 AI 平台（以下简称目标平台）及其工具链约束，围绕输入预处理、类别顺序、输出张量、后处理参数和运行时约束整理迁移检查流程，使训练权重、ONNX 中间表示、OM 离线模型和板端应用之间的关键约束能够被逐项检查。
    \item \textbf{面向安全视觉的部署行为对照实践。} 本文不只报告检测精度，还从转换、验证、敏感性、阈值行为和计时等角度分析模型迁移到板端后的实际表现，初步为电梯边缘视觉系统提供一种检查思路。
\end{enumerate}

这些内容分别对应后续章节。任务建模由第一章和第四章承接；迁移检查流程在第三章展开，覆盖输入、类别、输出、后处理和运行五类约束；部署行为对照则由第四章协议和第五章结果共同给出。

\section{研究内容组织逻辑与论文结构}

本文将研究重点放在“模型检测能力”和“边缘部署行为”两个层面。前者说明 YOLO 模型能否在验证集上识别行人与电动车，后者说明模型经过中间表示和离线模型格式进入边缘开发板后，能否继续保持输出含义清楚、结果可以追溯。基于这一思路，第二章介绍 YOLO 和 Hi3516DV500 板端部署基础，第三章给出兼容迁移方法，第四章定义数据、场景复杂性因素和评价协议，第五章基于验证结果分析设备端表现。这样安排后，模型训练、格式转换、设备端运行和实验分析能够串成连续链路。

本文正文共分为五章。第一章介绍研究背景、场景挑战、研究目标和论文结构。第二章回顾目标检测、YOLO 方法、评价指标和 Hi3516DV500 板端部署基础。第三章提出 YOLO 类双目标检测模型的边缘兼容迁移方法。第四章给出数据组织、场景复杂性因素和评测协议。第五章给出实验结果与分析，结论部分总结全文工作并讨论扩展方向。

\chapter{技术基础与相关工作}
\label{cha:background}

本章为后续方法设计铺垫必要的技术背景。本文关心的并不是单纯训练一个检测模型，而是让 YOLO 双目标检测模型迁移到 HongOU PI（Hi3516DV500）开发板后，输出仍然能够解释、能够复查。因此，本章围绕目标检测、YOLOv8、板端 AI 部署和迁移后的行为对照展开。

\section{目标检测基础}

目标检测要在同一张图像中给出目标位置和类别判断。它不同于只输出整图类别的图像分类，检测器最终需要返回边界框、类别标签和置信度分数。本文的类别只有 \textit{person} 与 \textit{ebike}，但电梯轿厢并不是一个容易场景。狭小空间、镜面反光、门框结构和近距离透视会让同一目标在不同帧中呈现不同尺度，边框回归和类别分数都会受到影响。

早期深度学习检测方法大多沿着“两阶段”路线展开：先生成候选区域，再完成分类和回归。R-CNN 利用卷积特征和候选区域分类提升检测精度\cite{girshick2014rcnn}，Fast R-CNN 将特征共享与端到端训练结合起来\cite{girshick2015fast}，Faster R-CNN 又用区域建议网络减少外部候选框依赖\cite{ren2017faster}。这一路线在定位和分类上较可靠，但推理步骤多，对边缘设备的算子和内存支持要求更高。SSD 把多尺度默认框与单阶段预测结合起来\cite{liu2016ssd}，YOLO 则把检测改写为端到端回归问题\cite{redmon2016yolo}，工程链路相对直接。

检测结果通常用 precision、recall、F1、IoU 和 mAP@0.5 描述。precision 关心预测框中有多少是真目标，recall 关心真实目标中有多少被找出来，F1 用来观察二者的折中；预测框与真实框的 IoU 达到阈值且类别一致时，才计为一次正确检测。本文采用 mAP@0.5 评价图像级检测能力，并在安全视觉分析中同时关注分类别召回率，因为电动车漏检比普通背景误检更直接影响预警效果。

mAP@0.5 可以概括一批图像上的平均表现，却不能替代板端行为检查。它无法说明应用程序是否按正确类别顺序读取输出张量，也不包含运行成功率、耗时口径和后处理阈值变化。平均指标还会遮住类别差异，例如 \textit{ebike} 与 \textit{person} 在遮挡形态、尺度变化和风险含义上并不相同。因此第五章同时给出总体指标、分类别指标、阈值后验重算和分段计时，让检测能力和设备端运行记录互相校验。

设某一类别的真正例、假正例和假负例分别为 TP、FP 和 FN，则 precision、recall 与 F1 可以写为
\begin{equation}
    P=\frac{TP}{TP+FP},\quad R=\frac{TP}{TP+FN},\quad F1=\frac{2PR}{P+R}.
\end{equation}
放到安全视觉语境中，precision、recall 和 F1 指向的风险并不一样。precision 偏低时，系统更容易产生无效提醒；recall 偏低时，真实风险可能没有被发现；F1 只是把二者放到同一个折中指标中观察。对于电梯场景，行人漏检会影响人员计数，电动车漏检则直接关系到入梯风险，两类召回率都需要单独查看。

检测框最终是否保留下来，还取决于置信度阈值和 NMS。置信度阈值控制低分候选框能否进入结果，NMS 阈值控制重叠候选框的抑制强度。阈值过高时，局部可见的电动车可能被漏掉；阈值过低时，门框、反光或大面积背景又更容易引入误检。因此，阈值不是普通工程参数，而是连接检测性能与安全输出的关键设置。

轿厢图像中的尺度变化并不只来自远近距离，还会受到近景截断和遮挡影响，检测器因此需要在不同分辨率特征之间取得平衡。FPN 将高层语义沿自顶向下路径传回浅层，并用横向连接补充空间细节\cite{lin2017fpn}；Focal Loss 处理的则是密集预测中前景样本少、易分类背景样本多的问题，通过降低简单样本权重让训练更关注难例\cite{lin2017focal}。在检测框架扩展方面，Mask R-CNN 将目标检测进一步延伸到实例分割\cite{he2017mask}，CenterNet 等 anchor-free 方法则直接预测中心点或关键点\cite{zhou2019centernet}。这些工作对本文的启发，不在于堆叠更复杂的检测范式，而在于提醒后续实验：候选框数量、类别分数和定位稳定性都可能受模型结构与后处理设置影响。本文因而选择 YOLO 作为可部署基础，并在部署端重点检查阈值、NMS 和应用层后处理行为。

电梯轿厢目标检测与开放场景检测存在明显差异。开放场景通常目标种类多、背景变化大，评价重点更多放在平均精度；电梯场景类别较少，但背景结构固定且反复出现，误报和漏报往往集中在少数高风险区域。表 \ref{tab:elevator_detection_diff} 概括了二者的差异。本文采用双目标检测而不是通用多类别检测，是为了让模型和评价协议聚焦于电梯安全视觉的核心需求——人员感知与电动车识别。

\begin{table}[htbp]
    \centering
    \caption{电梯场景检测与通用开放场景检测的差异}
    \label{tab:elevator_detection_diff}
    \begin{tabular}{p{2.8cm} p{4.2cm} p{4.6cm}}
        \toprule
        维度 & 通用开放场景 & 电梯轿厢场景 \\
        \midrule
        目标类别 & 类别多，目标分布复杂 & 类别少，但安全含义强 \\
        背景结构 & 背景变化大 & 门框、镜面和轿壁反复出现 \\
        目标尺度 & 尺度变化来自远近距离 & 近景截断和透视畸变明显 \\
        评价重点 & 平均精度和泛化能力 & 漏检、误检和连续输出稳定性 \\
        部署要求 & 可离线评测为主 & 需要端侧运行和结果可复查 \\
        \bottomrule
    \end{tabular}
\end{table}

\subsection{YOLO 系列检测器与 YOLOv8 特点}

YOLO 系列检测器将目标检测建模为单阶段预测问题，直接从输入图像生成边框、类别和置信度结果，具有结构紧凑、推理流程清晰和部署友好等特点\cite{redmon2017yolo9000}。与两阶段检测器相比，YOLO 系列更适合实时或准实时场景，这也是其在边缘视觉任务中被广泛采用的重要原因。

YOLO 的版本演化，一直在精度、速度和工程适配之间调整重心。早期的 YOLOv1 把检测统一为单次前向预测\cite{redmon2016yolo}，YOLO9000 和 YOLOv3 随后把多尺度训练、层次化类别学习和多尺度预测纳入模型设计\cite{redmon2017yolo9000}\cite{redmon2018yolov3}。到 YOLOv4、YOLOX 以后，数据增强、路径聚合、解耦头和 anchor-free 设计逐渐成为改进重点\cite{bochkovskiy2020yolov4}\cite{ge2021yolox}；YOLOv7、YOLOv6 也分别从训练技巧、工业部署和实时性角度推进单阶段检测器\cite{wang2022yolov7}\cite{li2022yolov6}。近年的 YOLOv9 将可编程梯度信息纳入训练设计\cite{wang2024yolov9}，YOLOv10 则把端到端无 NMS 检测作为探索方向\cite{wang2024yolov10}，之后的 YOLOv11、YOLO12 仍在延续这条工程化演进路线。本文并不把版本号本身作为比较重点；对于仅包含 \textit{person} 与 \textit{ebike} 的轿厢任务，更需要确认训练、导出、编译和板端执行链路能否稳定复查。

在本文使用的 YOLO 检测器中，图像先经过 Backbone 形成多层视觉特征，随后由 Neck 汇合浅层空间细节和高层语义信息，最后通过检测头输出框坐标与类别分数。这里强调结构划分，并不是为了重复网络模块名称，而是因为轿厢图像对特征融合比较敏感：近景人体常被画面边缘截断，电动车也可能只露出车把、车轮等局部结构。多尺度信息接合不足时，检测框容易在这些位置产生抖动或漏检。图 \ref{fig:yolo_pipeline} 给出了本文采用的 YOLO 双目标检测流程。

\begin{figure}[htbp]
    \centering
    \includegraphics[width=0.92\textwidth]{image/fig2-yolo-pipeline.png}
    \caption{YOLO 双目标检测流程}
    \label{fig:yolo_pipeline}
\end{figure}

本文选择 YOLOv8n，主要是出于工程链路和设备约束的共同考虑。YOLOv8 的训练、导出和推理生态相对成熟，YOLOv8n 的参数量和计算量较小，更适合 HongOU PI 这类边缘开发板。其 anchor-free 检测头和解耦结构也便于检查导出模型的输入输出。虽然 YOLOv9、YOLOv10 等新版本在部分基准上有提升，但本文的双目标任务更看重可部署、可复查和板端适配成本。模型只输出 \textit{person} 与 \textit{ebike} 两类分数；导出后，张量维度、坐标顺序和类别顺序仍必须和板端程序一致。若应用层沿用通用类别数或旧输出格式，即使权重本身没有问题，检测框也可能出现类别错位或坐标失真。

近年来，DETR 等端到端检测框架将目标检测建模为集合预测问题，减少了手工设计候选框和部分后处理步骤\cite{carion2020detr}。这类方法展示了目标检测结构的另一条发展路线，但在本文的边缘开发板场景中，模型部署成熟度、算子支持和实时处理成本同样重要。相比之下，YOLO 系列在工程链路中具有更成熟的导出、编译和板端适配经验，更适合作为本文可复查的边缘安全视觉实现基础。

YOLOv8 的主干、特征融合和检测头边界比较清楚，这有利于检查导出模型的输入输出。但结构清楚不代表部署过程会自动正确。训练框架中的预测结果通常已经完成坐标还原、类别筛选和 NMS；导出到板端后，程序拿到的可能只是未经封装的张量。应用层需要明确张量维度、类别数、分数含义和坐标尺度，否则训练阶段的指标无法和板端检测框放在同一口径下比较。

边框质量也需要统一口径。传统 IoU 衡量预测框与真实框的重叠程度，GIoU 将无重叠情况下的几何关系纳入计算\cite{rezatofighi2019giou}，DIoU 进一步考虑中心点距离\cite{zheng2020diou}。这些指标主要服务训练或离线评测，但它们提醒了一个板端问题：如果坐标解码、缩放或裁剪方式前后不一致，类别分数正常也不能保证 IoU 正常。第三章把边框解释放入接口契约，正是为了避免这类口径错位。

轻量化模型选择需要在精度和设备约束之间取舍。MobileNet、EfficientNet、EfficientDet 等研究讨论了特征提取效率、模型缩放和检测性能之间的关系\cite{howard2019mobilenetv3,tan2019efficientnet,tan2020efficientdet}。电梯场景类别数少，但这并不意味着模型表达能力可以降低要求；遮挡和局部可见目标仍然需要多尺度特征支持。另一方面，边缘开发板的算力和算子支持又限制了模型复杂度。本文采用轻量 YOLO 模型，正是出于检测能力和可部署性的折中。

轻量骨干网络研究为这一选择提供了背景。Inception 结构通过多分支卷积提升表达效率\cite{szegedy2015inception}，ResNet 用残差连接缓解深层网络训练退化\cite{he2016resnet}；MobileNet 系列通过深度可分离卷积和倒残差结构显著降低计算量\cite{howard2017mobilenets}\cite{sandler2018mobilenetv2}，SqueezeNet 以极小参数量实现接近 AlexNet 的精度\cite{iandola2016squeezenet}，ShuffleNet V2 总结了移动端网络设计的实际准则\cite{ma2018shufflenetv2}。MnasNet 和 MobileNetV3 进一步把硬件延迟纳入搜索目标\cite{tan2018mnasnet}\cite{howard2019mobilenetv3}，Once-for-All 则尝试通过一次训练适配多种部署约束\cite{cai2019onceforall}。这些研究共同说明，边缘检测不能只看理论 FLOPs，还需要考虑算子、内存访问和目标硬件的真实执行特性。

在本文任务中，YOLO 的错误往往能在轿厢结构中找到原因。门口的垂直边缘可能干扰边框判断，镜面反光会带来重复候选框，电动车只露出车把或车轮时，分数通常更接近阈值边界。模型结构、数据分布和后处理参数共同决定最终输出。后续实验保留分类别召回率和阈值扫描，目的就是把这些现象落到可以比较的指标上。

\section{边缘部署技术基础}

\subsection{Hi3516DV500 板端 AI 部署基础}

边缘智能强调在靠近数据源的位置完成推理，以降低带宽、隐私和响应时延压力\cite{shi2016edge}。本文使用的 HongOU PI（Hi3516DV500）并非仅作为通用开发板出现，而是作为面向电梯安全视觉的国产边缘 AI 实验平台。该平台搭载华为海思 Hi3516DV500 视觉 SoC，集成双核 ARM Cortex-A55 CPU、专用 NPU、图像处理与媒体输入输出能力，产品规格中给出的典型配置包括 2 GB DDR4、8 GB eMMC 和 2 TOPS INT8 NN 算力。与桌面 GPU 或服务器推理不同，Hi3516DV500 板端部署要求模型先从训练框架导出为中间表示，再通过 CANN/ATC 等工具编译为设备可执行的 OM 离线模型，并在板端应用中通过 SVP\_NPU 相关运行接口完成推理与输出解释。

因此，HongOU PI（Hi3516DV500）在本文中的意义不只是提供运行环境，而是把国产边缘视觉 SoC 的工具链约束纳入研究对象。其部署形态与昇腾生态中常见的“离线模型编译 + 板端运行时调用”工程范式相近，并体现了海思边缘 AI 平台在模型格式、算子支持、内存布局、交叉编译、后处理和运行时接口上的约束。本文选择该平台，是为了在真实板端链路中检验 YOLO 双目标检测模型从训练权重到 OM 模型再到应用输出的语义一致性。图 \ref{fig:ascend_deploy_pipeline} 展示了本文关注的部署链路。

\begin{figure}[htbp]
    \centering
    \includegraphics[width=\textwidth]{image/fig3-ascend-deploy.png}
    \caption{HongOU PI（Hi3516DV500）开发板上的 YOLO 部署链路}
    \label{fig:ascend_deploy_pipeline}
\end{figure}

这条链路从训练阶段的 YOLO 权重出发，经过 ONNX 导出、ATC 编译与优化，形成 OM 离线模型，再交给板端 NPU 推理。应用层还要继续完成输出解码、阈值筛选和 NMS。链路中的每一步都有可能改写最终结果：预处理差异会改变输入分布，张量读取顺序错误会改变类别语义，后处理阈值不同则会改变检测框数量。

在 Hi3516DV500 板端部署链路中，AIPP、输入尺寸、算子支持和内存布局都会影响模型执行。AIPP 若承担颜色空间转换或归一化操作，则训练阶段预处理必须与板端配置对应；若预处理在应用层完成，则应用层需要保证输入张量与模型期望一致。OM 模型编译完成后，其输入输出名称、形状和数据类型也需要被记录，否则后续调试很难判断错误来自模型本身还是应用层解释。对电梯安全视觉而言，这些约束会直接影响类别顺序、边框解码、候选框筛选和报警稳定性，本文因而把板端工具链和运行时接口作为部署一致性评价的一部分。

深度学习系统研究也表明，模型执行效率不仅由网络本身决定，还由编译器、运行时和硬件体系共同决定。TPU 的系统分析展示了专用加速器在数据中心推理和训练中的性能收益\cite{jouppi2017tpu}，TVM 则代表了通过端到端编译优化把模型映射到不同硬件后端的研究方向\cite{chen2018tvm}。边缘场景与数据中心不同，设备功耗、内存和算子支持更受限制，本文重点关注模型在 Hi3516DV500 板端完整执行后的输出解释和运行口径，而不是把通用加速器吞吐作为主要评价指标。

模型压缩和量化是边缘部署中的常见方法。整数推理友好量化、剪枝和知识蒸馏分别从数值表示、结构冗余和知识迁移角度降低模型部署成本\cite{jacob2018quantization,han2016deep,hinton2015distilling}。混合精度量化研究进一步表明，不同网络层对量化误差的敏感性并不相同，量化策略会影响最终精度和部署效率\cite{dong2019hawq}。不过，本文并不把模型压缩作为主要创新，而是关注场景化训练后的 YOLOv8n 模型在既定边缘平台上的兼容迁移与行为一致性。

量化研究进一步提示，端侧部署中的数值变化需要通过实测而不是假设处理。Krishnamoorthi 总结了卷积网络推理量化中的校准、缩放和整数计算问题\cite{krishnamoorthi2018quantizing}，Gholami 等对高效神经网络推理量化方法进行了系统综述\cite{gholami2021quantization}，AQD 则直接面向量化目标检测的精度问题\cite{chen2021aqd}。这些工作说明，检测模型量化后可能在分类分数、边框回归和候选框排序上产生差异。本文虽然不把量化算法作为贡献，但在实验中区分完整板端验证、阈值后验重算和分段计时，正是为了避免把部署后的数值行为简单等同于训练阶段模型。

\subsection{后处理、NMS 与部署参数敏感性}

目标检测的最终输出并不完全由神经网络前向结果决定。置信度阈值决定候选框是否进入后续统计，NMS 决定重叠框的保留关系，类别特定的清理规则则可能改变某一类目标的误报和漏报。Soft-NMS 表明，即使只改变传统 NMS 的抑制方式，也可能改善检测结果\cite{bodla2017softnms}。对于本文的电梯场景，NMS 的影响更具体：人和电动车经常在狭窄空间中重叠，电动车局部结构又可能生成多个相邻候选框，过强或过弱的抑制都可能改变最终安全输出。

本文不把阈值、NMS 和后处理策略简单看作“调参”。训练脚本与板端应用若采用不同置信度阈值，precision、recall 和预测框数量都会变化；NMS 阈值一改，拥挤重叠区域的候选框保留关系也会随之改变。应用层若增加类别特定清理规则，某些看似模型能力不足的问题，实际可能来自候选框处理。第五章中的阈值扫描、NMS 实验和后处理消融，正是为了把这几类影响拆开。

把后处理作为部署模块单独讨论，可以分清模型前向结果和应用层规则的边界。本文先报告端到端板端基准，再用参数敏感性和消融实验解释关键设计选择。这样写，是为了说明模型权重、阈值、NMS 与电动车清理规则如何共同决定最终输出。

边缘部署还要区分“能够运行”和“适合应用”。前者只说明模型可以编译、加载和推理；后者还要求输出含义清楚、失败样本有迹可循，后处理规则也要符合安全任务。本文在结果章同时报告成功样本数、batch 验证链路耗时、分段计时和阈值扫描，把板端可运行性拆成检测输出、验证链路和耗时口径几个可检查部分。

\section{相邻视觉任务与本文重点}

电梯轿厢内的安全视觉任务可以按观察尺度分成不同方向。轿厢拥挤程度通常要落到人数统计上，已有研究把边缘轻量检测器用于电梯人数统计系统\cite{lin2023elevatorcounting}。当问题转向目标在多帧中的持续状态时，DeepSORT、ByteTrack 等跟踪方法才成为主要工具\cite{wojke2017deepsort,zhang2022bytetrack}。姿态估计和视频识别还能给出人体关键点或动作模式\cite{cao2017openpose,feichtenhofer2019slowfast}，不过这些任务会引入视频级标注、时序模型和行为定义。本文的实验范围因此收束在 \textit{person} 与 \textit{ebike} 的板端检测，而不是把跟踪和动作识别一并纳入主线。

连续视频在本文中承担的是补充检查作用，用来观察单帧指标不容易暴露的框保持、闪烁和耗时波动。跟踪方法本身已有较成熟的处理路径：SORT 用卡尔曼滤波预测目标状态，再通过匈牙利匹配连接相邻帧检测框\cite{bewley2016sort}；DeepSORT 在此基础上加入外观特征，以改善跨帧身份关联\cite{wojke2017deepsort}；ByteTrack 则把部分低分检测框保留下来参与关联，减少轨迹断裂\cite{zhang2022bytetrack}。这些设计关注的是检测结果在时间轴上的延续关系，与本文的单帧双目标检测评价并不属于同一层级。因此，代表性视频片段只用于补充观察连续输出保持、闪烁和分段耗时，为板端部署一致性提供场景证据。

后续章节仍围绕图像级双目标检测和板端迁移展开；跟踪或动作识别可以作为后续扩展，但不放入本文的主要实验链路。

\subsection{部署一致性问题}

本文所说的部署一致性，指同一检测模型经过格式转换和设备迁移后，输入、输出和后处理仍处在可比较口径下。普通离线实验往往关注最终 mAP；安全视觉系统则不同，类别错位、边框漂移、阈值变化或运行时异常都可能改变预警结果。因此，这一问题不能只写成实验环境说明，而要进入方法设计。

在本文中，需要核对的内容包括类别、边框、阈值和运行记录。类别顺序要在训练、ONNX 导出、OM 编译和板端应用之间保持一致；坐标归一化、缩放和裁剪方式要能支撑同一套 IoU 判断；置信度阈值与 NMS 参数在不同阶段应表达相同含义；设备端还要记录成功样本数、失败样本数、batch 端到端耗时和分段计时。少了其中任何一项，板端指标都可能被误读。

这些检查项之间存在前后依赖。类别顺序错误时，precision 和 recall 的语义会失效；坐标还原错误时，模型可能被误判为定位能力不足；阈值口径不一致时，安全输出会受偶然参数影响；运行记录不足时，一次实验也难以代表真实设备行为。第二章先把这些问题说清楚，第三章再给出接口契约，实验章据此解释结果。

已有部署研究指出，深度模型从训练框架迁移到生产环境时，格式转换、算子支持和运行时差异会带来额外风险\cite{stacker2021deployment,chen2021aqd}。电梯场景会把这种风险放大：反光和门口结构使候选框分数更容易贴近阈值，局部可见电动车会让边框更敏感，拥挤重叠也会影响 NMS。后续方法章围绕接口契约和分阶段验证展开，就是为了把这些风险落到可检查的技术项上。

\section{相关工作小结}

本章的技术背景指向一个明确缺口：YOLO 系列适合作为电梯双目标检测基础，边缘 AI 平台也具备本地预警的硬件条件；但既有工作常把检测精度和边缘部署分开讨论，较少把训练权重、ONNX 中间表示、OM 离线模型和板端后处理放在同一条证据链里。本文以行人与电动车检测为任务载体，面向 HongOU PI（Hi3516DV500）开发板整理兼容迁移方法，并在后续章节用板端实测结果检验这条链路。

\chapter{YOLO 类双目标检测模型的边缘兼容迁移方法}
\label{cha:method}

本章给出论文的核心方法。这里不把模型训练和板端运行分开处理，而是把 YOLO 类双目标检测模型从训练权重迁移到目标平台、再进入设备端执行的过程视为一条需要保持语义一致的系统链路。围绕 HongOU PI（Hi3516DV500）开发板和 YOLOv8n 模型，模型训练、ONNX 导出、OM 编译和板端验证被组织为连续流程，使 \textit{person} 与 \textit{ebike} 两类目标的输入、类别、边框、阈值和运行状态都有可检查的约束。

\section{总体技术路线}

图 \ref{fig:migration_pipeline} 给出了方法总体流程。训练阶段得到的 YOLO 双目标检测权重先经过类别数与检测头适配检查，再导出为 ONNX 中间表示；随后通过 ATC 编译得到边缘设备可执行的 OM 模型；进入板端后，系统完成 NPU 推理、输出解析、阈值筛选和后处理，并与统一实验协议下的结果进行一致性分析。

\begin{figure}[htbp]
    \centering
    \includegraphics[width=\textwidth]{image/fig4-migration-method.png}
    \caption{YOLO 双目标检测模型的边缘兼容迁移流程}
    \label{fig:migration_pipeline}
\end{figure}

这条流程的关键不只在模型文件能否生成。模型层需要保证检测头输出与双目标任务一致，避免类别数或类别顺序错误；格式层要检查 PyTorch、ONNX 和 OM 之间的输入输出张量语义是否一致；应用层则要让板端后处理逻辑与训练框架中的统计口径保持可比较。YOLOv5、YOLOv8 等 YOLO 类检测器在部署时都会涉及输入预处理、类别维度、检测头输出、候选框解码、置信度筛选和 NMS 等环节，因此本章按照这些共性接口组织迁移方法，并在 YOLOv8n 板端链路中展开具体实现与实验分析。只有这些约束同时成立，边缘设备端的检测指标才有解释意义。

迁移检查在论文中分为语义、数值和运行三个层次。语义检查包括类别数量、类别顺序和输出张量维度；数值检查包括边框坐标、置信度范围和 NMS 后检测数量；运行检查包括模型加载、推理成功率、验证链路端到端耗时和关键分段耗时。这样的顺序先确认基础语义，再讨论 mAP 等检测指标，同时把 batch 外围链路耗时与模型执行阶段耗时分开解释，避免把不同运行口径混在一起。

\section{YOLOv8n 场景化训练与验证性能}

进入板端部署前，双目标检测模型先在 Ultralytics YOLOv8n 框架下完成场景化训练。训练环境为 x86\_64 工作站，主要硬件包括 Intel Xeon Gold 6226R CPU、约 125 GiB 系统内存和 4 张 NVIDIA GeForce RTX 3090 GPU；本次训练使用第 0 号 GPU。训练任务使用 \textit{person} 与 \textit{ebike} 两类标注，输入尺寸为 $640 \times 640$，batch size 为 16，训练 100 个 epoch，并以 COCO 预训练 YOLOv8n 权重作为初始化。该 YOLOv8n 模型对应后续 ONNX 导出、ATC 编译和 OM 板端部署的源权重，训练阶段验证结果也就成为分析板端指标回退幅度的参照。

图 \ref{fig:chap03_training_results} 给出了训练过程中的损失和验证集指标曲线。可以看到，训练后期 box loss、classification loss 和 DFL loss 继续保持下降或趋稳，验证集 precision、recall 和 mAP 指标整体维持在较高水平，未出现明显的训练发散现象。

\begin{figure}[htbp]
    \centering
    \includegraphics[width=\textwidth]{image/generated/fig_chap03_training_results.png}
    \caption{YOLOv8n 场景化训练曲线}
    \label{fig:chap03_training_results}
\end{figure}

第 100 个 epoch 的验证集 precision 为 0.9756，recall 为 0.96469，mAP@0.5 为 0.98832，mAP@0.5:0.95 为 0.82203。从这个结果看，场景化训练后的 YOLOv8n 模型在训练阶段已经具备较强的双目标检测能力。第五章进一步将训练阶段验证结果与 Hi3516DV500 板端 FP16/OM 结果对照，用于解释部署后指标的保留程度和回退幅度。

\begin{figure}[htbp]
    \centering
    \includegraphics[width=\textwidth]{image/generated/fig_chap03_training_val_predictions.jpg}
    \caption{YOLOv8n 场景化训练模型验证集预测示例}
    \label{fig:chap03_training_val_predictions}
\end{figure}

图 \ref{fig:chap03_training_val_predictions} 展示了训练阶段模型在验证集上的预测输出。该图来自 YOLOv8n 训练流程保存的验证预测网格，用于直观说明模型已经能够在电梯视角下同时输出 \textit{person} 与 \textit{ebike} 检测框；训练阶段能力的定量判断仍以验证集 precision、recall 和 mAP 指标为依据。

\section{迁移约束与一致性检查}

迁移链路在本文中被抽象为接口契约。这里的契约不是额外算法模块，而是贯穿训练、导出、编译和板端运行的一组约束，用来把容易隐含在代码和配置中的输入、类别、输出、后处理和运行时假设显式化。对于 YOLO 类检测器而言，这些契约分别落在输入图像语义、类别维度、检测头张量、候选框解码、置信度阈值、NMS 策略和板端运行状态等环节。只有这些约束在各阶段保持一致，板端 precision、recall、F1 和 mAP@0.5 才能对应到真实检测行为。

输入契约约束输入尺寸、颜色通道、归一化、缩放与填充方式。电梯图像常包含强反光和大面积背景结构，若训练阶段和板端的缩放、填充或颜色通道处理不一致，模型对局部电动车结构的响应会发生变化，进而表现为分数漂移或小目标漏检。因此，输入契约要求训练配置、ONNX 输入形状、OM 模型输入和板端预处理路径采用一致的图像语义。

类别契约约束类别数量、类别顺序和标签名称。论文采用 \textit{person} 与 \textit{ebike} 两类标签，类别集合较小，但类别解释必须在训练配置、导出模型、板端标签表和后处理代码中完全一致。若类别维度读取或标签顺序发生错位，检测框位置即使正确，报警语义也会失效。

输出契约约束检测头张量维度、坐标格式和置信度含义。训练框架通常自动完成坐标解码和结果封装，而板端程序需要显式读取输出张量并解释其维度。如果坐标中心点、宽高、左上右下角或置信度分数组合方式被错误理解，最终结果会表现为边框漂移、类别分数异常或检测数量变化。

后处理契约约束置信度阈值、NMS 阈值和类别处理方式。电梯场景中的反光、门框和局部可见电动车会生成接近阈值的候选框，阈值和 NMS 的小幅变化可能改变最终安全输出。因此，后处理契约要求训练阶段统计脚本和板端应用层在阈值含义、NMS 输入、类别过滤和电动车清理策略上保持可比较。

运行契约约束模型文件、板端二进制、运行模式、成功样本数、失败样本数、端到端处理耗时和分段计时。该契约用于区分模型执行阶段、单帧处理路径和 batch 验证链路，避免把外围链路耗时误解为 OM/NPU 模型执行耗时，也避免在样本未处理完时报告全量验证结论。

为了让契约不只停留在概念层面，迁移链路采用逐项检查方式。ONNX 导出后检查输入尺寸和输出张量形状是否符合双目标检测设置；OM 编译后检查板端程序读取的输出数量和张量维度是否与导出阶段一致；运行验证集时，再看成功样本数、分类别预测数和分类别 TP/FP/FN 是否处于合理范围。若某一项出现异常，优先回到对应契约定位问题，而不是直接调整阈值掩盖错误。

分阶段验证策略强调每一阶段只解决少数关键问题。训练阶段主要看模型是否具备基础检测能力，并保留权重文件与验证集指标；导出阶段检查输入输出形状和类别维度是否保留；编译阶段关注 OM 模型能否生成以及算子是否兼容；板端阶段再检查推理成功率、后处理结果、检测框、分类别指标和运行时信息。这样做可以降低排错复杂度，使板端结果异常时能够回溯问题来源。

下面对五类约束的关键一致性环节分别说明。

\subsection{输入预处理对齐}

输入预处理对齐要求训练阶段、导出阶段和板端使用相同的输入尺寸与图像变换语义。对于固定视角电梯场景，目标尺度分布相对集中，但近景人体和局部电动车会使边缘区域信息非常重要。因此，预处理不应仅追求图像尺寸匹配，还需要保证缩放比例、填充区域和坐标还原方式能够被后处理正确反推。

\subsection{类别语义对齐}

类别语义对齐是双目标任务迁移中最基础的契约。本文包含 \textit{person} 与 \textit{ebike} 两个类别，类别数较少反而要求类别解释更加清晰。训练配置、导出模型、板端标签表和后处理代码必须使用相同类别顺序；否则，即使检测框位置正确，最终报警语义也可能被反转。

在实验解释中，分类别 TP、FP、FN 统计也被用作类别契约的外部检查。若板端程序把类别维度读取错误，\textit{person} 与 \textit{ebike} 的预测数量、召回率和 mAP@0.5 会出现与场景经验明显不符的异常。因此，分类别指标不仅衡量检测性能，也用于判断迁移后的语义解释是否成立。

\subsection{输出张量语义对齐}

输出张量语义对齐要求板端程序明确知道 YOLO 检测头输出中每一维的含义。本文关注的双目标检测只有两个类别，但这反而要求类别维度更清晰：如果板端按错误类别数读取张量，类别分数的位置会发生偏移；如果按错误坐标格式解释边框，IoU 指标会下降。因此，输出张量检查是迁移流程中的必要步骤，它决定后续实验结果能否对应到真实检测行为。

\subsection{后处理一致性}

后处理一致性包括置信度筛选和 NMS 两部分。置信度阈值决定候选框进入最终结果的机会，NMS 决定重叠候选框的保留关系。电梯场景中的反光、门口结构和局部可见目标会产生大量边界候选框，因此阈值口径会直接影响安全输出。实验章通过阈值敏感性分析观察不同置信度阈值下 precision、recall、F1 和分类别召回率的变化。

后处理既可以在模型图中完成，也可以在应用层完成。若后处理前移到模型图中，板端应用更简单，但阈值和 NMS 策略调整不够灵活；若后处理保留在应用层，便于针对安全场景调整阈值，但应用层必须正确实现坐标解码和候选框抑制。本文选择应用层后处理思路，并通过契约方式约束其输入、阈值和输出含义。

\subsection{运行时可靠性}

运行时可靠性要求设备端不仅能够加载 OM 模型，还能够完成验证集推理并记录成功样本数、失败样本数、验证链路端到端耗时和关键分段耗时。对于安全视觉系统，单张图像推理成功不足以说明系统可用；只有在统一数据集上可靠执行，并能说明耗时指标属于哪一段链路，设备端结果才具备分析意义。

运行时可靠性还需要关注异常样本处理。若某些图像推理失败，实验报告应区分模型检测错误和运行失败；若大量样本触发兜底路径或超时，端到端耗时也需要与成功样本数一起解释。结果章会报告成功样本数、失败样本数、batch 图像验证链路平均单图处理耗时和分段计时结果，使精度指标能够与板端运行状态共同解释。

\section{一致性分析方法}

本文从类别语义、边框解释、阈值行为和运行时状态四个角度分析迁移后一致性。类别语义检查不同阶段是否输出相同类别含义；边框解释关注坐标还原、IoU 衡量和定位漂移设计；阈值行为用于观察置信度阈值变化对 precision、recall 和 F1 的影响；运行时状态则记录成功样本数、失败样本数、端到端处理耗时和分段计时。这些内容共同构成本文的分析框架。

具体到方法层面，类别一致性依靠类别顺序和分类别 TP/FP/FN 判断 \textit{person} 与 \textit{ebike} 语义是否可靠；边框解释借助 IoU、边框漂移设计和 NMS 后数量判断定位变化是否影响检测结果；阈值口径通过多阈值 P/R/F1 与 \textit{ebike} recall 判断安全输出是否依赖单一参数；运行时可靠性则通过成功样本数、失败样本数、端到端耗时和分段计时判断板端验证是否执行完毕，并区分 batch 外围链路与模型执行阶段。

跨阶段逐样本匹配是迁移后一致性分析的扩展形式。实验章重点报告板端验证集结果、分类别结果、阈值敏感性后验重算、NMS 敏感性、后处理消融、视频输出状态和分段计时结果；PyTorch、ONNX、OM 和板端四阶段的逐样本匹配可用于进一步细化分析类别、坐标和分数在迁移链路中的变化。这样的处理使实验结论与证据范围保持对应。

样本级匹配的基本思路是：对于同一张输入图像，先按类别过滤候选框，再按 IoU 从高到低匹配同一目标，并比较不同阶段的类别、坐标和分数差异。若类别一致但 IoU 明显下降，主要问题可能来自坐标还原或 NMS；若 IoU 接近但类别不同，问题更可能来自类别顺序或输出维度解释；若同一阶段检测数量明显变化，则需要检查阈值和候选框抑制逻辑。该方法使跨阶段一致性从抽象概念转化为可执行的比较过程。

当结果出现异常时，本文先判断异常表现属于类别、边框、阈值还是运行状态，再回到对应契约检查。若类别大量错位，优先检查类别顺序和输出张量维度；若边框整体偏移，优先检查缩放、填充和坐标解释；若检测数量突变，优先检查 conf、NMS 和类别清理策略；若时延异常升高，优先检查板端执行路径、输入输出复制和资源状态。这种排查顺序把部署问题转化为可定位的技术项，便于复用。

\section{本章小结}

本章整理了面向目标平台约束的 YOLO 类双目标检测模型兼容迁移流程。流程围绕迁移约束展开，把输入预处理、类别配置、输出张量、后处理参数和运行状态纳入同一检查框架。相较于单纯报告模型精度，论文更关注模型进入板端后的行为能否解释、结果能否追溯。

\chapter{数据集、场景复杂性与评测协议}
\label{cha:data_protocol}

\section{数据集与任务定义}

\subsection{任务与数据组织}

本文实验对象是电梯轿厢内的 \textit{person + ebike} 双目标检测。数据来自实际轿厢监控视频，抽帧后由人工使用 LabelImg 按 YOLO 格式标注，类别为 \textit{person} 和 \textit{ebike}。训练集约 3600 张标注图像，验证集固定为 720 张图像；遇到遮挡目标或局部可见目标时，标注采用可见区域外接矩形。图像验证集负责给出检测指标，代表性视频片段用于观察连续画面中的输出变化，板端运行记录则用于拆分阈值行为、运行状态和分段耗时。

实验设计围绕第一章的问题链展开，但不把指标堆成单纯列表。需要回答的核心问题包括：模型进入板端后能否跑完验证集，\textit{person} 与 \textit{ebike} 的分类别结果是否符合任务语义，置信度阈值、NMS 和应用层清理会不会改变安全输出，代表性视频片段中检测框是否连续，以及 batch 图像验证耗时、单帧核心处理耗时和模型执行阶段耗时能否区分清楚。第五章的结果组织按这些问题展开。

\subsection{场景复杂性因素}

电梯场景的困难之处来自封闭空间、固定视角和安全响应要求的叠加。轿厢内常见的高亮区域、门框结构、局部目标、近距离重叠、尺度变化和暗部区域都会改变目标轮廓、类别分数或候选框排序。这些因素在本文中作为误差解释维度，服务于阈值、NMS、后处理和视频稳定性实验中的现象分析。

本文采用的证据来自几类材料：720 张图像验证集、既有板端检测输出、NMS 参数实验、后处理消融、batch 图像复查、连续视频分段计时和代表性视频片段。图像验证集给出总体精度、分类别精度和 batch 验证耗时；阈值后验重算观察同一批输出在不同安全阈值下的变化；NMS 与后处理实验用来区分模型输出和部署规则对候选框的影响；分段计时拆开 batch 端到端耗时、单帧核心处理和模型执行阶段；视频片段用于观察连续帧中的漏检、闪烁和边框抖动。上述证据共同回答一个问题：同一 YOLO 模型迁移到 OM 并在板端运行后，检测行为是否仍能被解释和复查。

场景因素在本文中保持为误差解释维度。它们与图像级检测指标、板端运行记录一起使用，既避免论文主线从双目标检测发散开，也保留了电梯场景相对于通用检测任务的特殊性。

\section{评测协议设计}

\subsection{评测指标}

本文按用途组织指标，而不是简单分成若干并列类别。基础识别能力由 precision、recall、F1 和 mAP@0.5 描述；\textit{person} 与 \textit{ebike} 的 TP、FP、FN、recall 和 mAP@0.5 用来检查分类别语义；多阈值 F1、类别召回、NMS 参数下的 FP/FN 变化和后处理消融，用来观察部署参数对最终输出的影响。运行时部分则记录成功处理样本数、batch 图像验证链路端到端单图处理耗时和分段计时。跨阶段 IoU 漂移保留为更细粒度的方法设计，不作为本章主要实测结果。

\subsection{实验协议}

为保证结果可比较，图像级指标全部在同一验证集上统计，并统一保留三位小数。阈值实验只改变置信度阈值，其余参数保持不变；NMS 实验只改变 IoU 阈值；后处理消融只改变电动车清理策略；视频片段实验只改变平滑窗口；分段计时固定模型、二进制和阈值设置，分别记录 batch 端到端、单帧核心处理和模型执行阶段。本文并不追求实验数量堆叠，而是让每个实验只回答一个主要问题。

\subsection{场景复杂性与部署误差的关系}

场景复杂性会影响模型迁移后的结果解释。高亮区域和暗部区域常先影响类别置信度，在阈值扫描中表现为召回率波动；轿厢边界和门框结构更容易干扰边框回归与 NMS，进而拉低 IoU；近距离重叠则可能同时改变类别排序和候选框抑制关系，使不同阶段的输出数量不一致。

因此，实验分析需要把指标变化放回场景中理解。mAP@0.5 下降时，先判断问题来自类别分数还是边框定位；视频片段出现闪烁时，再结合阈值结果判断是否存在置信度边界问题。


\section{数据可追溯性与实验设计}

\subsection{数据集划分与可追溯性}

实验结果要能复查，数据划分必须先说明清楚。本文图像验证集固定为 720 张，用于报告板端总体指标和分类别指标。论文正文只交代验证集规模、标注类别和使用目的，不展开本地文件组织细节，以免把工程路径写成论文证据。

可追溯性主要体现在样本数和统计口径上。实验表格中的样本数需要与指标来源一致；\textit{person} 与 \textit{ebike} 的真实目标数、预测数、TP、FP、FN 要能对应到评测结果；成功样本数、失败样本数和端到端耗时则用于说明板端验证是否按预期执行。

视频片段采用检测保持率、闪烁窗口、连续无输出片段和分段计时等机器侧统计描述输出波动。这样可以让指标和实验材料相匹配。


\subsection{跨阶段一致性实验设计}

跨阶段一致性实验的理想形式是选取同一批输入样本，分别在训练阶段、ONNX 导出阶段、OM 编译阶段和板端部署阶段运行，再对输出结果进行逐样本匹配。匹配时先按照类别和 IoU 阈值确定同一目标，然后统计类别一致率、平均边框 IoU、置信度差异和最终检测数量差异。该实验能够直接衡量模型迁移过程是否改变检测行为。

在实际执行时，应避免同时改变多个变量。输入图像、预处理方式、置信度阈值和 NMS 阈值需要固定，唯一变化的是推理阶段。若发现输出差异，再根据第三章的误差定位表检查类别契约、坐标契约或后处理契约。这样，跨阶段实验不仅给出结果，还能帮助定位问题来源。

本文已经积累了板端全量验证、参数敏感性、后处理消融、视频连续输出和分段计时证据。训练阶段、ONNX 导出阶段、OM 编译阶段和板端部署阶段的逐样本统一对比可在此基础上进一步展开，本节给出其分析方式，用于衔接后续更细粒度的迁移分析。


\subsection{实验结果呈现原则}

为了保证实验结果可信，本文区分实测结果、后验重算结果和方法性评价设计。已有完整证据支撑的内容，如 720 张图像的板端验证结果、NMS 参数敏感性、后处理消融、代表性视频稳定性和分段计时结果，可以写入实验结果并进行结论分析；基于既有检测输出重新筛选得到的阈值敏感性结果，作为补充实验处理；跨阶段逐样本对齐则保留为进一步细粒度分析的评价方法。

上述做法与实验研究对证据分层的基本要求一致。工程实践中存在日志、脚本、视频和人工观察等多种材料，只有区分证据强度，才能避免把经验观察混同为定量结论。本文用图像验证集计算检测指标，用参数敏感性和后处理消融解释部署规则影响，用视频片段辅助观察稳定性，再用场景复杂性描述解释误差来源。

经过这样的划分，实验结果、后验重算和方法设计各有作用，部署一致性分析可以在实测证据基础上逐层展开。

\section{本章小结}

本章给出了本文的数据组织方式、场景复杂性因素、评测指标体系和实验协议。场景复杂性描述用于误差解释，分析重点放在 YOLO 双目标检测模型迁移到 HongOU PI（Hi3516DV500）开发板后的部署行为是否对齐。

\chapter{实验结果与分析}
\label{cha:experiments}

本章按照第四章定义的实验协议，分析 YOLO 双目标检测模型在边缘设备端的表现。实验并不只报告一个离线精度数值，而是先用官方基线、YOLOv8n 场景化训练结果和板端 FP16/OM 结果建立参照系，再分别考察全量验证集检测能力、分类别语义、置信度阈值、NMS 参数、应用层后处理、代表性视频片段和分段计时口径。这样安排后，板端结果可以和训练阶段能力、部署规则以及运行时状态放在同一证据链中解释。

\section{实验设置与评价范围}

本文实验使用同一轻量 YOLO 双目标检测模型，类别为 \textit{person} 与 \textit{ebike}。图像验证集包含 720 张电梯轿厢图像，评价指标包括 precision、recall、F1、mAP@0.5、分类别指标和板端运行状态。板端推理使用搭载华为海思 Hi3516DV500 视觉 SoC 的 HongOU PI 开发板完成，模型以 OM 离线模型形式运行，后处理包含置信度筛选和 NMS。

表 \ref{tab:board_platform_specs} 给出了本文实验使用的板端平台与部署工具链。该信息用于限定本文所有板端指标的硬件和运行口径，后续检测指标、参数敏感性、视频稳定性和分段计时结果均基于同一类部署环境获得。

\begin{table}[htbp]
    \centering
    \caption{板端实验平台与部署工具链}
    \label{tab:board_platform_specs}
    \begin{tabular}{p{3.0cm} p{7.6cm}}
        \toprule
        项目 & 配置 \\
        \midrule
        开发板 & HongOU PI（Hi3516DV500） \\
        平台定位 & 国产边缘 AI 视觉 SoC 开发平台 \\
        核心 SoC & 华为海思 Hi3516DV500 视觉 SoC \\
        CPU & 双核 ARM Cortex-A55 @ 850 MHz \\
        DDR & 2 GB DDR4 \\
        eMMC & 8 GB \\
        NPU & 2 TOPS INT8 NN 算力 \\
        模型格式 & OM 离线模型 \\
        部署工具链 & ATC / CANN / SVP\_NPU \\
        交叉编译目标 & AArch64 GNU/glibc 工具链 \\
        \bottomrule
    \end{tabular}
\end{table}

图 \ref{fig:board_photo} 展示了本文实验使用的 HongOU PI（Hi3516DV500）开发板实物。该开发板体积紧凑，适合嵌入电梯控制柜或轿厢顶部安装，板载接口支持摄像头视频输入和网络数据传输。

\begin{figure}[htbp]
    \centering
    \includegraphics[width=0.55\textwidth]{image/fig_board_photo.png}
    \caption{HongOU PI（Hi3516DV500）开发板实物}
    \label{fig:board_photo}
\end{figure}

实验前先确认板端验证环境、模型文件、应用程序和数据集均处于就绪状态。全量验证集已完成板端实测，阈值敏感性分析采用既有板端检测输出进行后验筛选重算；NMS 和后处理消融采用同一板端 campaign 协议，保证只有一个参数变化；分段计时实验则使用新增计时埋点对小样本 batch 图像验证和代表性视频片段进行补充复查。

阈值后验重算的依据是板端检测输出中保留了每个候选框的类别、置信度和坐标，同时每张图像具有对应真实标注。因此，在原始候选阈值低于待分析阈值的前提下，可以通过重新筛选候选框并按 IoU@0.5 匹配真实框，得到不同置信度阈值下的 precision、recall 和 F1。该做法不改变模型推理结果，只改变结果解释阈值，适合用于分析安全输出对阈值选择的敏感程度。

本章评价问题与证据形式保持一一对应。RQ1 使用 720 张图像完整验证总体指标回答板端完整验证是否可执行且具备检测能力；RQ2 使用分类别 TP、FP、FN、recall 和 mAP@0.5 检查两类目标是否保持正确语义；RQ3 使用 0.20--0.40 置信度阈值后验重算分析安全输出稳定性；RQ4 和 RQ5 分别使用固定 score 后改变 NMS 阈值、改变电动车清理策略的板端 campaign；RQ6 使用代表性视频片段和平滑窗口对比观察连续输出；RQ7 使用 50 张图像 batch 复查与连续视频分段计时解释运行时指标。这样的组织方式用于保证每个实验结论都能追溯到具体数据类型，避免把不同口径的指标混为一谈。

\section{板端检测能力验证}

\subsection{板端完整验证集结果}

表 \ref{tab:board_overall_results} 给出了 720 张图像验证集上的总体结果。所有样本均在板端完成处理，成功处理样本数为 720，失败样本数为 0。总体 precision 为 0.901，recall 为 0.923，F1 为 0.912，mAP@0.5 为 0.910。板端 batch 图像验证链路端到端平均单图处理耗时约为 1119.510 ms。

\begin{table}[htbp]
    \centering
    \caption{边缘设备端总体检测结果}
    \label{tab:board_overall_results}
    \begin{tabular}{l c}
        \toprule
        指标 & 数值 \\
        \midrule
        验证图像数 & 720 \\
        成功样本数 & 720 \\
        失败样本数 & 0 \\
        precision & 0.901 \\
        recall & 0.923 \\
        F1 & 0.912 \\
        mAP@0.5 & 0.910 \\
        batch 端到端单图处理耗时/ms & 1119.510 \\
        \bottomrule
    \end{tabular}
\end{table}

这组数据至少说明两点。第一，YOLOv8n 场景化训练模型经过 ONNX 导出、ATC 编译和 OM 部署后，能够在边缘设备上跑完 720 张验证图像。第二，板端输出可以继续计算 precision、recall、F1 和 mAP@0.5，说明后续分类别分析、阈值扫描和运行时拆解有共同的数据基础。

表 \ref{tab:board_overall_results} 中的 1119.510 ms 是 batch 图像验证链路的端到端单图处理耗时，包含图像输入、媒体链路、预处理、模型执行、后处理和结果写出等环节。它描述的是验证路径整体成本，不能等同于 OM 模型执行阶段耗时，也不能直接代表连续视频链路的帧间处理时间。后续分段计时实验会专门拆开这一差异。

\subsection{训练阶段与板端同集对比}

为了避免将板端指标孤立理解，本文把第五章结果放入同一任务证据链中解释。证据链包括官方 YOLOv8n COCO-person 基线、基于 RTX 3090 GPU 完成训练的 YOLOv8n 场景化模型，以及 Hi3516DV500 板端 FP16/OM 完整验证结果。前者用于说明通用预训练模型直接迁移到电梯俯视场景时的不足，中间结果作为训练阶段能力参照，板端结果则用于衡量部署后检测能力是否被保留。三者共同回答的问题不是“本文指标在所有公开数据集上排名如何”，而是在同一电梯验证集和同一安全任务背景下，模型从通用基线、场景训练到板端部署的效果变化是否具有解释性。

表 \ref{tab:stage_comparison_evidence} 给出了三类证据的来源和可比性边界。官方 YOLOv8n 基线采用已修正的 COCO-person 逐图 IoU 匹配评估口径，precision 为 0.743，recall 为 0.728，F1 为 0.736。该基线只评价 \textit{person} 类别，不能与双类别 mAP@0.5 直接等价，但可以作为通用 COCO person 能力在电梯俯视场景中的低基线。YOLOv8n 场景化训练模型第 100 个 epoch 的 precision 为 0.976，recall 为 0.965，mAP@0.5 为 0.988；板端 FP16/OM 完整验证的 precision 为 0.901，recall 为 0.923，F1 为 0.912，mAP@0.5 为 0.910。

\begin{table}[htbp]
    \centering
    \caption{训练阶段与板端同集对比证据}
    \label{tab:stage_comparison_evidence}
    \begingroup
    \zihao{5}
    \setlength{\tabcolsep}{4pt}
    \renewcommand{\arraystretch}{1.25}
    \begin{tabular}{@{}p{2.4cm}p{3.0cm}p{3.2cm}p{4.2cm}@{}}
        \toprule
        阶段 & 环境与格式 & 核心指标 & 可比性说明 \\
        \midrule
        COCO-person 基线 & \makecell[l]{官方 YOLOv8n\\逐图评估} & \makecell[l]{P=0.743\\R=0.728\\F1=0.736} & 只评价 person 类别，用于说明通用模型不足 \\
        YOLOv8n 场景化训练 & \makecell[l]{PyTorch/Ultralytics\\RTX 3090 GPU} & \makecell[l]{P=0.976\\R=0.965\\mAP@0.5=0.988} & 训练阶段双目标检测参照 \\
        Hi3516DV500 板端 & \makecell[l]{OM 离线模型\\FP16/板端应用} & \makecell[l]{P=0.901\\R=0.923\\F1=0.912\\mAP@0.5=0.910} & 部署后完整 720 张验证集实测 \\
        \bottomrule
    \end{tabular}
    \endgroup
\end{table}

图 \ref{fig:chap05_stage_comparison} 展示了阶段间指标关系。左侧使用 precision、recall 和 F1 比较三类证据，可以看到 YOLOv8n 场景化训练模型相比官方 COCO-person 基线有明显提升，说明电梯轿厢中的俯视视角、遮挡和目标形态差异需要场景化训练支撑。右侧比较训练阶段验证结果与板端 FP16/OM 结果的 mAP@0.5，板端结果相比训练阶段下降约 0.078。这一差异来自模型导出、OM 编译、板端预处理、输出解码和后处理链路的共同影响；不过板端 mAP@0.5 仍达到 0.910，主要检测能力在部署后得到了保留。

\begin{figure}[htbp]
    \centering
    \includegraphics[width=\textwidth]{image/generated/fig_chap05_stage_comparison.pdf}
    \caption{官方基线、YOLOv8n 训练阶段与板端 FP16/OM 结果对比}
    \label{fig:chap05_stage_comparison}
\end{figure}

后续阈值敏感性、NMS 敏感性、后处理消融和分段计时实验，并不是对一个孤立板端数值的简单补充，而是在“训练阶段能力已经建立、板端部署后存在可量化回退”的前提下，继续解释部署链路中哪些因素会影响最终安全输出。

\subsection{分类别输出分析}

表 \ref{tab:board_class_results} 给出了 \textit{person} 与 \textit{ebike} 两类目标的检测结果。电动车类别的 mAP@0.5 为 0.952，召回率为 0.961，高于行人类别的 0.869 和 0.888。这说明在当前验证集上，模型对电动车目标具有较强识别能力，而行人类别受到遮挡、截断和多人重叠等因素影响更明显。

\begin{table}[htbp]
    \centering
    \caption{边缘设备端分类别检测结果}
    \label{tab:board_class_results}
    \begin{tabular}{lrrrrrrr}
        \toprule
        类别 & GT & Pred & TP & FP & FN & recall & mAP@0.5 \\
        \midrule
        person & 776 & 794 & 689 & 105 & 87 & 0.888 & 0.869 \\
        ebike & 720 & 739 & 692 & 47 & 28 & 0.961 & 0.952 \\
        \bottomrule
    \end{tabular}
\end{table}

图 \ref{fig:chap05_detection_metrics} 对总体指标和分类别指标进行了可视化。整体四项指标均处于 0.90 左右，说明板端输出具备基本检测能力；分类别结果中，电动车类别的召回率和 mAP@0.5 均高于行人类别，符合电动车结构特征更清晰、行人更容易受近景截断和多人重叠影响的场景观察。

\begin{figure}[htbp]
    \centering
    \includegraphics[width=\textwidth]{image/generated/fig_chap05_detection_metrics.pdf}
    \caption{边缘设备端总体指标与分类别指标}
    \label{fig:chap05_detection_metrics}
\end{figure}

为补充分类别指标的直观解释，图 \ref{fig:chap05_board_val_examples} 给出了从 720 张完整验证集板端输出中按固定规则抽取的检测示例。图中底图来自验证集原图，检测框和置信度来自 HongOU PI（Hi3516DV500）开发板生成的检测记录，并按板端记录的帧尺寸恢复到图像坐标。示例筛选条件为样本运行成功、同时包含 \textit{person} 与 \textit{ebike}、两类均产生 TP 且无 FP/FN；在满足上述条件的样本中，优先选择检测置信度较高、目标边界不贴边且可视化较清晰的样本。该图用于展示板端输出形态，完整性能判断仍以表 \ref{tab:board_class_results} 和图 \ref{fig:chap05_detection_metrics} 的全量统计为准。

\begin{figure}[htbp]
    \centering
    \includegraphics[width=\textwidth]{image/generated/fig_chap05_board_val_examples.png}
    \caption{Hi3516DV500 板端检测输出示例}
    \label{fig:chap05_board_val_examples}
\end{figure}

从部署行为对照的角度看，分类别统计还可以反向验证类别契约是否成立。如果类别顺序在导出或板端解析过程中发生错位，\textit{person} 与 \textit{ebike} 的 TP、FP、FN 分布会出现明显异常，难以呈现当前这种与场景经验一致的分类别结果。因此，分类别指标不仅用于衡量检测性能，也是部署语义检查的一部分。

\section{部署参数敏感性分析}

\subsection{阈值敏感性分析}

为分析置信度阈值对安全输出的影响，本文在既有板端检测输出基础上进行后验筛选重算。由于原始板端运行使用的候选阈值低于本节扫描阈值，本文可以在不重新推理的条件下，分别保留置信度不低于 0.20、0.25、0.30、0.35 和 0.40 的检测框，并按 IoU@0.5 重新统计 precision、recall 和 F1。表 \ref{tab:threshold_sensitivity_results} 给出了结果。

\begin{table}[htbp]
    \centering
    \caption{置信度阈值敏感性后验重算结果}
    \label{tab:threshold_sensitivity_results}
    \begin{tabular}{cccccc}
        \toprule
        conf & precision & recall & F1 & ebike recall & person recall \\
        \midrule
        0.20 & 0.911 & 0.919 & 0.915 & 0.957 & 0.884 \\
        0.25 & 0.921 & 0.916 & 0.919 & 0.957 & 0.878 \\
        0.30 & 0.925 & 0.912 & 0.919 & 0.956 & 0.872 \\
        0.35 & 0.935 & 0.909 & 0.922 & 0.956 & 0.866 \\
        0.40 & 0.942 & 0.903 & 0.922 & 0.953 & 0.857 \\
        \bottomrule
    \end{tabular}
\end{table}

随着置信度阈值从 0.20 提高到 0.40，总体 precision 从 0.911 提升到 0.942，recall 从 0.919 下降到 0.903，F1 保持在 0.915--0.922 区间。电动车召回率在五个阈值下保持在 0.953 以上，说明在当前验证集和板端输出条件下，电动车检测对这一阈值区间不太敏感。行人召回率随阈值提高从 0.884 降至 0.857，说明行人类别中存在更多接近阈值边界的候选框。

图 \ref{fig:chap05_threshold_sensitivity} 展示了各指标随阈值变化的趋势。precision 随阈值提高而上升，recall 缓慢下降，F1 保持相对平稳；分类别召回率中，\textit{ebike} 曲线较平，\textit{person} 曲线下降更明显。阈值选择不能只依据总体 F1，还需要考虑不同类别在安全任务中的风险权重。

\begin{figure}[htbp]
    \centering
    \includegraphics[width=\textwidth]{image/generated/fig_chap05_threshold_sensitivity.pdf}
    \caption{置信度阈值敏感性曲线}
    \label{fig:chap05_threshold_sensitivity}
\end{figure}

阈值敏感性后验重算并不等同于重新训练模型，也不改变板端推理输出。它的作用是分析同一批板端输出在不同安全阈值下的波动情况，因此可以增强本文关于当前验证集内阈值行为的证据链。

\subsection{NMS 参数敏感性}

NMS 参数敏感性实验关注的是重叠候选框如何被保留。实验固定模型、验证集、置信度阈值和后处理策略，只改变 NMS 阈值，取值为 0.35、0.45 和 0.55。阈值较低时，相邻候选框会被更强地抑制，重复框可能减少，但遮挡区域也更容易漏检；阈值较高时，局部目标的候选框保留更多，误报和重复框也可能随之增加。

这组实验不同于阈值后验重算，需要重新运行板端推理，因为 NMS 已经在板端应用层改变最终输出。三组实验使用同一 OM 模型、同一板端二进制和同一验证集切分，均处理 720 张图像，失败样本数为 0。

结果显示，NMS 从 0.35 提高到 0.55 时，预测框总数从 1526 增至 1542，FP 从 146 增至 160，FN 从 116 降至 114。行人类别 F1 在三组设置下均为 0.878，变化很小；电动车 F1 从 0.952 降至 0.944，而电动车召回率从 0.960 小幅升至 0.963。可以看出，提高 NMS 阈值会多保留一些电动车候选框，漏检略少，但误报也增加，因此 F1 略有下降。

\begin{figure}[htbp]
    \centering
    \includegraphics[width=\textwidth]{image/generated/fig_chap05_nms_sensitivity.pdf}
    \caption{NMS 参数敏感性结果}
    \label{fig:chap05_nms_sensitivity}
\end{figure}

从安全场景角度看，NMS=0.45 是较平衡的设置：其 mAP@0.5 为 0.910，整体 FP/FN 分别为 152 和 115，与 0.55 的 mAP@0.5 基本相当，但误报更少；与 0.35 相比，0.45 保留了略多候选框并取得更高 mAP@0.5。本文采用的 NMS 设置并不是孤立参数，而是在召回、误报和重复框保留之间取得折中。

\subsection{后处理策略消融}

应用层后处理消融用于回答一个容易混淆的问题：板端输出变化究竟来自模型权重，还是来自模型之后的候选框清理逻辑。后处理消融中，电动车误检清理策略被显式暴露为 full、safe 和 off 三种模式。其中，full 表示完整启用误检清理，safe 表示只启用较保守的重复框和明显条带清理，off 表示关闭电动车类别特定清理。默认 auto 模式保持当前板端行为，即 batch 评价路径保留更原始的 safe 输出，视频 file 路径使用更面向展示的 full 输出。

该消融不是训练消融，也不表示本文提出了新的检测网络结构。它的作用是把板端应用层规则从“隐含工程细节”转化为可以复查的部署模块。三组实验均完成 720 张验证图像，失败样本数为 0，产物分别对应 cleanup\_full、cleanup\_safe 和 cleanup\_off campaign。

full 策略将电动车 FP 从 safe 的 47 降至 38，但电动车 FN 从 28 增至 169，电动车 recall 从 0.961 降至 0.765，F1 从 0.949 降至 0.842。也就是说，完整清理策略确实减少了部分电动车误报，但代价是过度抑制真实电动车候选框。相比之下，safe 策略保留较高电动车召回率；off 策略进一步将电动车 recall 提高到 0.963，但 FP 增至 55，F1 降至 0.944。

\begin{figure}[htbp]
    \centering
    \includegraphics[width=\textwidth]{image/generated/fig_chap05_postprocess_ablation.pdf}
    \caption{后处理策略消融结果}
    \label{fig:chap05_postprocess_ablation}
\end{figure}

这一组结果提醒本文不能只讨论模型权重，应用层后处理策略同样会改变最终输出。对于图像验证链路，safe 策略在误报控制和召回保持之间更平衡，因此更适合作为板端 batch 分析的默认口径；full 策略可用于展示链路中的强清理，但不宜直接作为图像级检测指标的默认设置。正文评价以完整验证集统计为主，定性样例和视频片段用于补充说明输出形态与连续性。

\section{运行时评价}

\subsection{视频稳定性观察}

除图像验证集外，本文还使用代表性视频片段观察连续场景下的输出波动。机器侧统计显示，在一个约 14.58 s、437 帧的代表性片段中，系统能够记录原始检测输出、筛选后输出、分段计时和运行资源信息。视频片段评价采用检测保持率、闪烁窗口、连续无输出片段和分段计时等机器侧指标，用于观察连续帧中的检测保持、闪烁、漏检现象和运行时状态。

视频稳定性实验固定模型、score 和 NMS，只改变平滑窗口，重点比较窗口 1 与窗口 5。前者接近无平滑或最小平滑，能够暴露单帧检测分数波动；后者代表默认连续输出策略，更接近真实展示链路。两组实验均处理 437 帧，产物分别为 video\_smooth1 和 video\_smooth5。

从逐帧输出看，两组实验的检测保持率均为 0.982，均出现 3 个 person\_flash 闪烁窗口，闪烁帧数均为 5，连续无输出片段数均为 2，最长无输出片段为 5 帧。窗口 1 与窗口 5 在该代表性片段上的输出统计基本一致，说明当前片段中主要波动不是简单通过增大平滑窗口即可消除的单帧噪声。两组视频路径的单帧处理耗时均值分别为 11.811 ms 和 11.951 ms，模型执行耗时均值分别为 8.359 ms 和 8.365 ms，均低于 30 fps 对应的 33.3 ms 帧间隔。

\begin{figure}[htbp]
    \centering
    \includegraphics[width=\textwidth]{image/generated/fig_chap05_video_stability.pdf}
    \caption{代表性视频片段稳定性结果}
    \label{fig:chap05_video_stability}
\end{figure}

视频片段分析表明，安全视觉系统的最终输出不仅取决于单帧检测框，还取决于后处理策略和连续帧稳定性。即使图像级指标较高，若连续视频中目标框频繁出现和消失，仍会影响预警体验。因此，本文将视频分析作为图像验证集的补充。本节的“检测保持率”和“连续无输出片段”描述连续输出状态，图像级检测精度仍由 720 张验证集结果承担。

\subsection{分段计时分析}

为避免把 batch 图像验证链路端到端耗时误解为 OM 模型执行阶段耗时，本文补充了板端分段计时实验。该实验在应用程序中记录单帧处理、模型执行、预处理、后处理和渲染等分段字段。其中，端到端字段表示 batch 图像验证链路单图耗时，单帧处理字段表示核心处理路径耗时，模型执行字段表示包围 OM/NPU 模型执行调用的工程计时。

表 \ref{tab:timing_breakdown_results} 给出了本轮分段计时结果。50 张图像 batch 复查中，端到端耗时均值为 1123.300 ms，与 720 张完整验证集中的 1119.510 ms 处于同一数量级；但同一实验中的单帧处理耗时均值为 21.459 ms，模型执行耗时均值为 8.208 ms。代表性视频片段中，单帧处理耗时均值为 11.870 ms，模型执行耗时均值为 8.407 ms。

\begin{table}[htbp]
    \centering
    \caption{板端分段计时结果}
    \label{tab:timing_breakdown_results}
    \begin{tabular}{lrrrr}
        \toprule
        运行模式 & 样本规模 & 端到端耗时 & 单帧处理耗时 & 模型执行耗时 \\
        \midrule
        batch 图像验证 & 50 张图像 & 1123.300 & 21.459 & 8.208 \\
        连续视频处理 & 437 帧 & -- & 11.870 & 8.407 \\
        \bottomrule
    \end{tabular}
\end{table}

图 \ref{fig:chap05_timing_breakdown} 使用对数坐标展示上述差异，图中虚线表示 30 fps 视频约 33.3 ms 的帧间隔。可以看到，batch 端到端耗时与单帧核心处理、模型执行阶段不在同一数量级。这说明 1119.510 ms 指标主要用于描述 batch 图像验证路径的端到端成本，而不应解释为模型执行阶段耗时。代表性视频路径下的单帧核心处理耗时低于 33.3 ms，说明该按源帧节奏回放条件下的处理计时与连续帧输入节奏相容；这一结论属于运行时分段证据，用于补充图像验证集指标。

\begin{figure}[htbp]
    \centering
    \includegraphics[width=\textwidth]{image/generated/fig_chap05_timing_breakdown.pdf}
    \caption{板端分段计时结果对比（对数坐标）}
    \label{fig:chap05_timing_breakdown}
\end{figure}

分段计时结果的适用范围由运行模式决定。50 张图像 batch 复查用于解释运行时口径，并与 720 张完整验证集检测指标形成互补；代表性视频片段说明连续处理路径和运行时耗时；模型执行耗时则是当前模型、二进制、阈值和板端环境下的工程观测值，用于佐证本文关于部署行为是否对齐的分析。

\section{综合讨论}

\subsection{部署一致性综合分析}

把训练阶段和板端结果放在一起看，官方 COCO-person 基线不足以支撑电梯俯视场景检测，场景化训练后的 YOLOv8n 明显提高了同集检测能力。模型进入 Hi3516DV500 板端后，mAP@0.5 相比训练阶段回退约 0.078，但仍达到 0.910。这个数值不是训练指标的复制，而是经过 ONNX 导出、ATC 编译、OM 加载和应用层后处理之后得到的实测结果。

720 张验证图像全部处理成功，说明当前板端链路至少没有出现大规模运行失败。分类别统计也能正常给出 \textit{person} 与 \textit{ebike} 的 TP、FP、FN，未观察到类别顺序错位或输出张量无法解读的迹象。这并不意味着迁移过程没有任何损失，而是说明输入、类别、输出和运行记录在本轮验证范围内可以对上。

阈值结果进一步暴露了两类目标的差异。电动车召回率在 0.20--0.35 区间内变化较小，说明当前验证集中的多数电动车候选框分数离阈值边界较远；行人召回率下降更明显，近景遮挡和多人重叠仍会产生低置信度候选框。这个差异与第四章的场景因素相符，也解释了为什么行人类别在板端回退中更敏感。

NMS 实验给出的变化幅度较小，但方向比较清楚。0.35、0.45 和 0.55 三组设置均能处理 720 张图像；阈值升高后，电动车候选框数量和 FP 略增，FN 略降，说明 NMS 主要影响电动车相邻候选框的保留关系。综合 mAP@0.5、FP 和 FN，0.45 更适合作为本文主要实验设置。

后处理消融则提醒本文不能把所有输出变化都归因于模型权重。完整清理策略减少了电动车 FP，却使电动车 FN 增加，recall 降至 0.765；safe 与 off 保留了更高召回率，其中 safe 相比 off 少产生 8 个电动车 FP。也就是说，应用层清理规则本身就是部署链路中的变量，需要和权重效果分开解释。

视频片段和分段计时补充了图像验证集无法回答的问题。窗口 1 和窗口 5 的检测保持率、闪烁窗口数和最长无输出段一致，说明该片段中的输出波动不能简单归因于平滑窗口设置。batch 图像验证链路中的 1119.510 ms 可以说明验证路径可复查，但不能解释为 OM 模型执行本身约需 1.1 s；分段计时给出的模型执行阶段耗时，才对应更细的板端运行观察。

因此，本章结论来自几类证据的相互印证：板端全量验证说明链路可执行，分类别统计检查类别语义，阈值与 NMS 实验观察参数影响，后处理消融划清应用层规则的作用，视频片段和分段计时补足连续输出与运行口径。PyTorch、ONNX、OM 和板端四阶段的逐样本匹配仍可作为后续更细粒度的边框漂移分析路径。

\subsection{场景复杂性与误差解释}

第四章列出的场景因素，可以解释分类别结果中的差异。行人误差主要集中在近景截断、多人重叠和门口区域干扰上。人体靠近摄像头或只出现局部结构时，预测框与真实框容易重合不足；多人重叠时，NMS 可能抑制部分候选框，带来漏检或定位偏差。电动车整体召回率较高，但只露出车把、车轮或车身边缘的样本仍然困难，置信度也更容易贴近阈值边界。

阈值结果与这一判断相互印证。电动车召回率在 0.20--0.35 区间内变化较小，说明多数电动车候选框分数余量较充足；随着阈值升高，行人召回率下降更明显，说明该类别中有更多低置信度但可能正确的候选框。因此，设定置信度阈值时，不能只盯着总体 F1，还要把不同类别在安全任务中的风险权重纳入考虑。

场景因素的作用，是帮助判断误差来源。如果某类复杂样本在训练阶段和板端阶段都表现困难，问题更可能来自数据分布和模型表达；如果同一复杂样本训练阶段正确而板端异常，就需要优先检查输入、输出或后处理契约。本文已有的板端验证、阈值重算、NMS 实验、后处理消融、视频片段观察和分段计时，为后续跨阶段分析提供了基础。

\subsection{结果讨论与扩展方向}

综合这些结果，本文方法的有效性主要体现在检测能力保留和板端行为可复查两点。官方基线、YOLOv8n 场景化训练结果和板端 FP16/OM 结果说明，场景化训练能够提升电梯双目标检测能力，模型部署后仍保留主要检测效果。板端 720 张验证图像全部处理成功，也说明从训练权重到边缘设备执行的迁移链路可以跑通。分类别、阈值、NMS、后处理消融、视频连续输出观察和分段计时结果进一步说明，论文不仅报告检测指标，也对安全阈值、应用层规则、连续输出状态和运行时口径进行了拆解。

围绕进一步应用，后续工作可继续从运行链路、视频分析和跨阶段匹配三个方向推进。batch 图像验证链路端到端耗时较高，可进一步优化图像读取、数据搬运和外围调度，使实验验证链路与连续视频链路在运行口径上更加接近。代表性视频片段可扩展为带逐帧标注的视频集，从而把连续输出的一致程度与视频级检测精度结合起来。跨阶段逐样本对比则可进一步量化每个目标从训练阶段到板端部署阶段的边框漂移，使兼容迁移分析更加细粒度。

\section{本章小结}

本章报告了 YOLO 双目标检测模型的训练阶段与板端同集对比、边缘设备端 720 张图像验证结果、分类别结果、阈值后验重算、NMS 实验、后处理消融、视频片段观察和分段计时。场景化训练后的 YOLOv8n 在训练阶段 mAP@0.5 为 0.988，板端 FP16/OM 在 720 张验证图像上取得 0.910 的 mAP@0.5，其中行人类别为 0.869、电动车类别为 0.952。多组阈值和 NMS 设置下，两类目标的召回率没有出现剧烈波动。后处理消融显示，完整清理策略会明显牺牲电动车召回率，safe 策略更适合作为图像验证链路默认口径；视频实验补充说明了连续输出状态和分段耗时。上述结果支撑了本文关于兼容迁移和板端行为检查的主要结论。

\chapter{结论}
\label{cha:conclusion}

\section*{1. 论文工作总结}
\addcontentsline{toc}{section}{1. 论文工作总结}

本文围绕电梯轿厢中行人与电动车的双目标检测，探讨了 YOLO 类模型从训练环境迁移到边缘设备后的行为变化。研究的出发点是：离线训练精度并不能直接代表板端实际表现，模型格式转换、输出解释和后处理策略都会引入额外的不确定性。

针对这一问题，本文的主要工作可归纳为三个层面。在任务层面，将电梯安全视觉的核心需求——人员感知与电动车识别——统一为双目标检测框架，并系统梳理了轿厢场景中影响检测精度的遮挡、反光和近距离透视等因素。在迁移层面，将训练权重到板端推理的转换过程组织为五类约束检查，使每一步格式变换中的输入尺寸、类别顺序、坐标编码和阈值语义都可被逐项核实。在评价层面，不再依赖单一的离线 mAP，而是从分类别指标、阈值敏感性、NMS 参数影响、后处理消融、分段计时和视频连续输出等多角度交叉验证板端行为。

实验结果揭示了若干值得关注的现象。板端 mAP@0.5 相比训练阶段回退约 0.078，其中行人类别的下降幅度明显大于电动车类别（0.869 vs 0.952），这与轿厢近景遮挡导致行人候选框置信度偏低有关。阈值敏感性实验表明，总体 F1 在较宽的置信度区间内保持稳定，说明板端模型的安全输出不依赖精确的单一阈值调优。后处理消融则进一步揭示，应用层清理策略对电动车召回率的影响远大于模型本身——完整清理策略使电动车 recall 骤降至 0.765，这意味着部署一致性评价必须将后处理模块从模型权重效果中显式剥离。

总体而言，本文工作表明，边缘安全视觉系统的评价不应止步于“模型能跑、精度达标”，而应覆盖从格式转换到后处理策略的完整部署链路。本文整理的分阶段检查方法和多维评价思路，可为同类电梯或封闭空间的边缘视觉部署提供参考。

\section*{2. 工作展望}
\addcontentsline{toc}{section}{2. 工作展望}

围绕进一步应用，后续可从四个方向推进。第一，扩充电梯轿厢数据来源，覆盖更多楼宇类型、摄像机安装角度、光照条件和人车共现形态，使双目标检测模型在更丰富的真实场景中得到验证。第二，在更多 YOLO 变体和边缘芯片平台上复用本文的迁移检查流程，比较不同模型与硬件组合的部署行为，进一步分析方法在同类海思/昇腾边缘 AI 平台及工具链中的适用性。第三，开展训练、ONNX、OM 和板端四阶段的逐样本匹配，量化边框、类别分数和候选框数量在迁移过程中的变化，使部署一致性评价更加细粒度。第四，扩充带逐帧标注的视频评价数据，将连续输出稳定性与视频级检测精度结合起来，并在同一框架下继续拓展人数统计、跌倒检测等电梯安全视觉任务。

\bibliographystyle{gbt7714-numerical}
\bibliography{main.bib}
%%
% 致谢
%%

\chapter{致谢}

这项本科毕业设计从开题、系统实现、设备验证到论文写作，几乎每一步都得到过老师和同学们的帮助。在论文完成之际，我想向在课题推进过程中给予指导、支持和鼓励的师长、同门、家人和伴侣表达感谢。

衷心感谢我的导师钟竞辉老师在课题选题、研究收敛与论文组织方面给予的持续指导。老师在研究问题凝练、技术路线选择以及学术表达规范等方面提出了许多具体建议，也让我在课题推进过程中逐渐理解工程实现、实验验证与论文写作之间的关系。

也感谢项目实践过程中提供帮助的同门。电梯边缘视觉系统的训练、导出、量化和设备端运行，需要跨越训练环境、模型转换环境与边缘设备等多个环节，许多问题是在一次次交流和排查中逐步弄清楚的。

同时，感谢家人和伴侣在学习与毕业设计期间给予的理解、陪伴与鼓励。正是因为有这些支持，我才能把精力持续投入到模型实验、系统验证与论文写作中。

再次感谢所有给予我帮助、支持与鼓励的人。没有这些指导、交流和陪伴，本次毕业设计和论文写作都不会像现在这样完成。

\end{document}
